% Part: computability
% Chapter: computability-theory
% Section: computable-sets

\documentclass[../../../include/open-logic-section]{subfiles}

\begin{document}

\olfileid{cmp}{thy}{cps}
\olsection{محاسبه کېدونکي سټونه}

د محاسبه کېدنې مفهوم له محاسبه کېدونکو تابعو څخه محاسبه کېدونکو
سټونو ته غځولے شو:

\begin{defn}
  فرض کړئ $S$ د طبيعي عددونو يو سټ دے۔ $S$ هله او يوازې هله
  \emph{محاسبه کېدونکے} دے چې ځانګړونکې تابع~$\Char{S}$ ئې محاسبه
  کېدونکې وي۔ په بله وينا، $S$ هله او يوازې هله محاسبه کېدونکے دے
  چې تابع
\[
\Char{S}(x) =
\begin{cases}
1 & \text{که $x \in S$ وي} \\
0 & \text{که داسې نۀ وي}
\end{cases}
\]
محاسبه کېدونکې وي۔ همدارنګه، اړيکه
$R(x_0, \dots, x_{k-1})$ هله او يوازې هله محاسبه کېدونکې ده چې
ځانګړونکې تابع ئې محاسبه کېدونکې وي۔

محاسبه کېدونکي سټونه او اړيکې \emph{د پرېکړې وړ} هم بلل کېږي۔
\end{defn}

\begin{explain}
پام وکړئ چې اوس د محاسبه کېدنې څو مفهومونه لرو: د جزوي تابعو، تابعو
او سټونو دپاره۔ دا سره ګډ نۀ کړئ!{} \iftag{TMs}{هغه ټيورينګ ماشين چې
  جزوي تابع محاسبه کوي، د تابع په هغو ننوتونو کښې ئې وت راګرځوي چې
  تابع پرې تعريف شوې وي؛ هغه ټيورينګ ماشين چې يو سټ محاسبه کوي، تر
  دې پرېکړې وروسته چې ننوت په سټ کښې دے که نه، يا~$1$ يا~$0$
  راګرځوي۔}{}
\end{explain}

\end{document}
